%mac
\subsubsection{GIA TỐC}
%\paragraph{Gia tốc} 
\Binhluan{Gia tốc là đại lượng vật lí đặc trưng cho sự biến đổi nhanh hay chậm của vận tốc.
\begin{align*}
\boder{\vec a=\dfrac{\Delta \vec v}{\Delta t}} 
\end{align*}
%\paragraph{Vector gia tốc} 
%\begin{align*}
%\boder{\vec{a}=\dfrac{\vec{v}-\vec{v}_0}{t-t_0}=\dfrac{\Delta \vec{v}}{\Delta t}}
%\end{align*}
%\begin{itemize}
%	\item[$\bullet$] Gốc tại vật chuyển động.
%	\item[$\bullet$] Phương không đổi theo phương quỹ đạo.
%	\item[$\bullet$] Chiều không đổi.
%\end{itemize}
\begin{center}
\begin{tabular}{|c|c|}
\hline 
$a.v > 0$ & $a.v < 0$ \\ 
\hline 
$\vec{a} \uparrow \uparrow \vec{v}$ & $\vec{a} \uparrow \downarrow \vec{v}$ \\ 
\hline 
vật chuyển động nhanh dần& vật chuyển động chậm dần\\ 
\hline 
\end{tabular} 
\end{center}}{$a=\dfrac{\Delta v}{\Delta t}$ được gọi là gia tốc trung bình.\\
Đơn vị của gia tốc là $m/s^2$}
\subsubsection{CHUYỂN ĐỘNG THẲNG BIẾN ĐỔI ĐỀU}
\newghichu{Độ lớn của vận tốc tức thời hoặc tăng đều hoặc giảm đều theo thời gian.\\
Gia tốc không đổi theo thời gian}
\begin{itemize}
\item Nếu tốc độ tức thời tăng đều thì đó là chuyển động nhanh dần đều. 
\item Nếu tốc độ tức thời giảm đều thì đó là chuyển động chậm dần đều.
\end{itemize}
\Binhluan{\Check \textbf{Công thức}
\begin{align*}\boder{v=v_0+at};\quad\boder{d=v_0t+\dfrac{1}{2}at^2}\end{align*}
\begin{align*}
\boder{v^2-v_0^2=2ad}
\end{align*}}{Khi sử dụng các công thức trên ta quy ước chọn chiều dương sao cho $v>0$. Khi đó: \\
$\bullet$ Nhanh dần đều: $a>0$\\
$\bullet$ Chậm dần đều: $a<0$}
\Note{Nếu ta chọn chiều dương tuỳ ý thì biểu diễn các $\vec v,\ \vec a$ và xét dấu chúng.
\begin{itemize}
\item Chuyển động thẳng nhanh dần đều: $a.v>0$.
\item Chuyển động thẳng chậm dần đều: $a.v<0$. 
\end{itemize}
}
\subsubsection{ĐỒ THỊ VẬN TỐC - THỜI GIAN CỦA CHUYỂN ĐỘNG THẲNG BIẾN ĐỔI ĐỀU}
\begin{itemize}
\item Vì $v=at+v_0\Leftrightarrow y=ax+b$ nên đồ thị có dạng đường thẳng.
\item Điểm xuất phát của đồ thị có hệ số góc là $a=\tan\alpha$. 
\begin{itemize}
\item Đồ thị đi lên nếu $a > 0$.
\item Đồ thị đi xuống nếu $a < 0$.
\end{itemize}
\centerline{\begin{tikzpicture}[scale=1,thick,>=stealth',draw=Mapcolor]
\tkzDefPoints{0/0/o, 0/3.5/x, 5/0/t, 0/0.5/x0}
\tkzDrawSegments[->,blue](o,x o,t)
\foreach  \x  in  {1.5,3,4.5}
\foreach  \y  in  {1,1.5,...,3}
{\draw[dotted,thin]  (\x,0)  --  (\x,3);
\draw[dotted,thin]  (0,\y)  --  (4.5,\y);}
\draw [Mapcolor, line width = 1.2pt](x0)--(4.5,3)coordinate(x1);
\draw [green!70!black, line width = 1.2pt,dashed](x0)--(4.5,0.5)coordinate(x2);
\tkzLabelPoint[right](x){$v$}
\tkzLabelPoint[above](t){$t$}
\tkzLabelPoint[left](x0){$v_0$}
\tkzLabelPoint[left](o){$O$}
\tkzDrawPoints[size=3,fill=white](o,x0)
\node at (2.25,-0.5){$a>0$};
\tkzMarkAngles[size=0.7cm](x2,x0,x1)
\tkzLabelAngles[pos=1](x2,x0,x1){$\alpha$}
\begin{scope}[shift={(7,0)}]
\tkzDefPoints{0/0/o, 0/3.5/x, 5/0/t, 0/2.5/x0}
\tkzDrawSegments[->,blue](o,x o,t)
\foreach  \x  in  {1.5,3,4.5}
\foreach  \y  in  {1,1.5,2,3}
{\draw[dotted,thin]  (\x,0)  --  (\x,3);
\draw[dotted,thin]  (0,\y)  --  (4.5,\y);}
\draw [Mapcolor, line width = 1.2pt](x0)--(4.5,0)coordinate(x1);
\draw [green!70!black, line width = 1.2pt,dashed](x0)--(4.5,2.5)coordinate(x2);
\tkzLabelPoint[right](x){$v$}
\tkzLabelPoint[above](t){$t$}
\tkzLabelPoint[left](x0){$v_0$}
\tkzLabelPoint[left](o){$O$}
\tkzDrawPoints[size=3,fill=white](o,x0)
\node at (2.25,-0.5){$a<0$};
\tkzMarkAngles[size=0.7cm](x1,x0,x2)
\tkzLabelAngles[pos=1](x1,x0,x2){$\alpha$}
\end{scope}
\end{tikzpicture}}
\end{itemize}